\begin{tomtat}
\begin{itemize}
  \item Doshi-Velez và Kim nói rằng đánh giá theo chức năng, cách đánh giá
  không có người tham gia, thích hợp nhất khi đối tượng đem đo đã được kiểm
  chứng từ trước bằng một cách đánh giá đắt hơn. Mọi chỉ số trung thực của
  chương~\ref{ch:do-do-trung-thuc} đều nằm ở cách đánh giá ấy, nên cả Phần~IV
  đang mang một món nợ kiểm chứng. Hình~\ref{fig:ch12-hai-nhanh-danh-gia} tách
  hai nhánh mà một lời giải thích đi qua và đánh dấu bằng nét chấm hai thứ
  đang thiếu: phép đối chiếu giữa hai nhánh, và quan hệ giữa lời giải thích
  với mô hình mà không nhánh nào đo.
  \item Cuộc kiểm đếm của Suh và cộng sự rà 18\,254 bài XAI trong cơ sở dữ liệu
  Scopus và cho một cái phễu bốn mức: 253 bài có từ khóa về thử nghiệm với
  người, 237 bài đúng chủ đề, 128 bài thật sự kiểm chứng với người, tức 0.70\%
  của mức đầu. Trong 237 bài đã tự tuyên bố có khả năng giải thích cho người,
  128 bài đem lời tuyên bố ra thử; bài báo phát biểu tỉ lệ ấy là 54\% ở một chỗ
  và 56\% ở chỗ khác, nên sách in phép đếm.
  \item Con số ấy không phải con số về độ trung thực. Các chữ
  \enquote{faithfulness}, \enquote{ground truth} và \enquote{attribution} không
  xuất hiện lần nào trong bài. Cái được đếm là việc kiểm chứng lời tuyên bố về
  khả năng giải thích cho người, tức phía tính hợp lý của cặp khái niệm mà
  mục~\ref{sec:ch01-do-trung-thuc-va-tinh-hop-ly} tách ra. Và 128 là chặn trên
  chứ không phải chặn dưới, vì phép chấm chỉ hỏi một thí nghiệm hành vi có tồn
  tại hay không: bài báo nêu ba chỗ mà một thí nghiệm có thể rỗng, là không
  biết ai được hỏi, hỏi sai đối tượng, và hỏi những câu mà người trả lời không
  kiểm tra được hoặc những câu đặt phương pháp cạnh một mốc so sánh đã hỏng
  sẵn. Lập luận trục của bài là lập luận tín hiệu: một lời tuyên bố về khả năng
  giải thích chỉ được xác nhận khi có xác nhận rằng tín hiệu đã được nhận và
  hiểu đúng.
  \item Nghiên cứu của Liedeker và cộng sự chạy đúng phép đối chiếu ấy, nhưng
  trên lời giải thích phản thực chứ không trên attribution đặc trưng, nên mọi
  con số của nó dừng ở ranh giới mà mục~\ref{sec:ch09-doi-tuong-khac-nhau} đã
  dựng cho bài báo neo. 167 người chấm 85 lời giải thích trên ba bộ dữ liệu
  bảng theo năm chiều, gộp thành một điểm duy nhất; cả năm chiều đều nằm ở phía
  tính hợp lý.
  \item Kết quả: gộp mọi lời giải thích thì chỉ Trust Score đạt mức ý nghĩa với
  điểm gộp, $r = 0.307$, còn lại đều dưới 0.1. Tách theo bộ dữ liệu thì các
  tương quan đổi cả dấu, độ lệch chuẩn trung bình của chúng giữa ba bộ là 0.31,
  và đó là con số nặng hơn cả, vì một chỉ số đổi dấu theo lĩnh vực thì không
  hiệu chỉnh được. Trên cả 127 tổ hợp con của bảy chỉ số, cả năm lớp mô hình dự
  báo đều cho $R^2$ trung bình âm, và thêm chỉ số quá ba tới bốn cái thì làm dự
  báo xấu đi.
  \item Cái đi qua ranh giới không phải con số mà là tiền lệ: phép đối chiếu
  giữa chỉ số tự động và cảm nhận của người là làm được, và kết quả của nó
  không hiển nhiên. Bài báo còn thuật lại rằng bộ đánh giá M4 cho thấy các chỉ
  số trung thực của attribution đặc trưng chỉ tương quan yếu với nhau và có thể
  cho ra những bảng xếp hạng mâu thuẫn; sách thuật qua lời bài báo vì không đọc
  M4. Đó là phép thử đồng thuận, không phải phép đối chiếu với ground truth,
  nên phát biểu bỏ ngỏ của
  mục~\ref{sec:ch09-doi-tuong-khac-nhau} không đổi.
\end{itemize}
\end{tomtat}

\cauhoi
\begin{cauhoids}
  \item Mục~\ref{sec:ch12-hai-nhanh-danh-gia} lập luận rằng dù hỏi bao nhiêu
  người thì câu trả lời của họ cũng không trực tiếp cho biết lời giải thích có
  phản ánh đúng mô hình hay không. Hãy nêu một thiết kế thí nghiệm có người mà
  kết quả của nó vẫn ràng buộc được độ trung thực, rồi chỉ ra thiết kế ấy phải
  giả định thêm điều gì về nhiệm vụ giao cho người tham gia.
  \item Cuộc kiểm đếm loại bài rà tổng quan, bài hội thảo nhỏ và bài không
  tiếng Anh ngay trong truy vấn. Với từng loại trong ba loại, hãy cho biết việc
  loại nó làm con số 0.70\% lệch lên hay lệch xuống, và vì sao.
  \item Trong bốn tình huống bị chấm 0 mà mục~\ref{sec:ch12-mot-cuoc-kiem-dem}
  liệt kê, tình huống thứ tư là bài so kết quả của mình với nhãn do người gán.
  Hãy giải thích vì sao tình huống ấy vẫn nằm trọn ở nhánh trên của
  hình~\ref{fig:ch12-hai-nhanh-danh-gia}, rồi cho biết bộ dữ liệu neo của
  chương~\ref{ch:chi-so-khong-do-do-trung-thuc} có rơi vào cùng tình huống hay
  không.
  \item Mục~\ref{sec:ch12-ket-qua-va-ranh-gioi} lập luận rằng độ lệch 0.31 giữa
  các bộ dữ liệu nặng hơn hệ số 0.307. Hãy dựng lại lập luận ấy bằng lời của
  mình, rồi nêu một trường hợp mà kết luận ngược lại: một chỉ số tương quan yếu
  nhưng ổn định vẫn tệ hơn một chỉ số tương quan mạnh nhưng đổi dấu.
  \item Completeness trong bảy chỉ số của bài báo được tính bằng SHAP. Hãy chỉ
  ra chuỗi giả định mà một con số Completeness phải dựa vào, tính từ tiên đề
  của chương~\ref{ch:shap} trở đi, rồi cho biết chuỗi ấy đứt ở đâu nếu
  attribution SHAP không trung thực.
  \item Kết quả về tổ hợp cho thấy thêm chỉ số vào quá ba tới bốn cái thì dự
  báo xấu đi. Hãy nêu hai cách giải thích khác nhau cho hiện tượng đó, một cách
  quy về bản chất của các chỉ số và một cách quy về cỡ mẫu, rồi cho biết dữ
  liệu nào phân biệt được hai cách giải thích ấy.
  \item Đọc \enquote{Fewer Than 1\% of Explainable AI Papers Validate
  Explainability with Humans} của Suh và cộng sự~\autocite{p24humangap}, mục
  phương pháp và bảng 1. Hãy sửa truy vấn Scopus in trong bảng ấy thành một
  truy vấn đếm số bài kiểm chứng độ trung thực thay vì khả năng giải thích cho
  người, rồi cho biết truy vấn mới gặp khó khăn gì mà truy vấn gốc không gặp.
\end{cauhoids}
